\documentclass[11pt]{article}

\usepackage[margin=1in]{geometry}
\usepackage{amsmath,amssymb,mathtools,bm}
\usepackage{booktabs,longtable,array}
\usepackage{enumitem}
\usepackage{microtype}
\usepackage{setspace}
\usepackage{xcolor}
\usepackage{hyperref}
\usepackage{cleveref}
\usepackage{threeparttable}
\usepackage{pdflscape}

\hypersetup{
  colorlinks=true,
  linkcolor=blue!55!black,
  citecolor=blue!55!black,
  urlcolor=blue!55!black
}

\setlength{\parindent}{0pt}
\setlength{\parskip}{0.55em}
\onehalfspacing
\setlist[itemize]{topsep=0.2em,itemsep=0.15em,leftmargin=1.4em}
\setlist[enumerate]{topsep=0.2em,itemsep=0.2em,leftmargin=1.6em}

\newcommand{\E}{\mathbb{E}}
\newcommand{\Var}{\operatorname{Var}}
\newcommand{\Cov}{\operatorname{Cov}}
\newcommand{\Normal}{\mathcal{N}}
\newcommand{\HSA}{\mathrm{HSA}}
\newcommand{\CES}{\mathrm{CES}}
\newcommand{\LCP}{\mathrm{LCP\text{-}HSA}}
\newcommand{\LHSA}{\mathrm{L\text{-}HSA}}
\newcommand{\LU}{\mathrm{L\text{-}U}}

\title{\textbf{Data, Model, and Estimation Design}\\
\large Competition and Inflation under HSA Demand: A Nine-Cell Identification-First Local State-Space Design}
\author{}
\date{Final Research Design Blueprint}

\begin{document}
\maketitle

\begin{abstract}
This document specifies a pre-declared nine-cell, identification-first semi-structural design for studying whether an empirical business-demography competition coordinate changes inflation dynamics through a competition-dependent marginal-cost or slack slope and a short-run market-participation channel. The nine cells combine three quarterly inflation/expectation blocks---PPI with GDP-deflator expectations, CPI with CPI expectations, and core CPI with CPI expectations---with three activity measures---inverse markup, a Beveridge--Nelson output gap, and a negative unemployment gap.

The cells are not given symmetric structural interpretations. The PPI--inverse-markup cell is the unique quantitative HSA cell and is used to evaluate the local cross-equation restriction
\[
\kappa_1=\zeta_0\theta_0.
\]
The remaining eight cells are pre-declared mechanism-robustness or transportability exercises and do not receive the quantitative HSA equality. Annual firm counts anchor the empirical competition coordinate, while quarterly business-demography indicators provide additional information about its stationary competition cycle. Full-joint and modular-cut analyses are reported side by side. Identification diagnostics are co-primary with model evidence.

The flexible-price/steady-state reference markup is estimated jointly rather than imposed in a preliminary calibration. The same reference markup determines the inverse-markup marginal-cost deviation and the HSA reference elasticity, eliminating sample-mean normalization from the theory-linked cell. Bayes factors measure integrated within-cell model evidence, time-series-respecting predictive scores evaluate predictive performance, and posterior predictive checks diagnose absolute fit. None of these evidence measures overrides weak-identification diagnostics.
\end{abstract}

\tableofcontents
\newpage

% ================================================================
\section{Objective and Scope}
% ================================================================

The design studies whether product-market competition changes inflation dynamics through two empirically distinct channels:
\begin{enumerate}
\item a competition-dependent marginal-cost or slack slope; and
\item a short-run market-participation channel whose loading may itself vary with the competitive environment.
\end{enumerate}

The analysis is deliberately semi-structural. The business-demography system is used to construct a latent empirical competition coordinate, while HSA theory supplies a local cross-equation restriction only in the empirical cell for which the inflation and marginal-cost measures are sufficiently close to the producer-pricing objects in the theory.

The empirical analysis is organized around a pre-declared \(3\times3\) matrix of nine data cells. The nine cells are not treated symmetrically in structural interpretation. One cell is the primary quantitative HSA cell; the remaining eight are mechanism-robustness or transportability exercises.

The design distinguishes:
\begin{itemize}
\item quantitative HSA consistency;
\item robustness of the competition mechanisms; and
\item transportability of those mechanisms across alternative inflation and slack measures.
\end{itemize}

No result from a secondary cell is used to redefine the primary HSA restriction ex post.

% ================================================================
\section{Research Questions}
% ================================================================

\subsection{RQ1: Competition-Dependent Slope}

Does the slow competitive environment change the response of inflation to marginal cost or macroeconomic slack?

The empirical object is
\begin{equation}
\kappa_1
=
\left.
\frac{\partial}{\partial \bar q}
\left(
\frac{\partial \pi}{\partial x}
\right)
\right|_{\bar q=0}.
\end{equation}

\subsection{RQ2: Short-Run Participation Channel}

Conditional on the maintained timing, measurement, and shock-orthogonality assumptions, is the short-run competition cycle associated with inflation through a distinct participation-channel loading?

The empirical object is
\begin{equation}
\theta_0
=
-
\left.
\frac{\partial\pi}{\partial\hat q}
\right|_{\bar q=0}.
\end{equation}

This coefficient is not interpreted as a nonparametrically identified causal effect.

\subsection{RQ3: Local HSA Restriction}

In the theory-near PPI/inverse-markup cell, are the slope and participation channels quantitatively linked as HSA predicts?

The local restriction is
\begin{equation}
\boxed{
\kappa_1=\zeta_0\theta_0.
}
\label{eq:local_hsa_restriction}
\end{equation}

For the unrestricted model define
\begin{equation}
\delta_{HSA}
=
\kappa_1-\zeta_0\theta_0.
\label{eq:delta_hsa}
\end{equation}

This restriction is tested only in the primary quantitative HSA cell.

\subsection{RQ4: State Dependence of the Participation Loading}

Does the participation-channel loading vary with the slow competitive environment?

Write
\begin{equation}
\theta(\bar q)
=
\theta_0+\gamma\bar q+O(\bar q^2).
\end{equation}

The primary empirical interpretation of \(\gamma\) is the local state dependence of the empirical participation-channel loading.

A direct interpretation of \(\gamma\) as a derivative of structural pass-through additionally requires that the mapping between the theoretical effective-variety coordinate and the empirical competition coordinate have locally constant scale at the reference environment.

Thus \(\gamma=0\) means local constancy of the empirical participation loading. It is interpreted as locally constant structural pass-through only under this additional proxy-mapping condition.

\subsection{RQ5: Transportability}

Are the signs, magnitudes, and state dependence of the competition mechanisms stable across alternative inflation indexes and alternative marginal-cost/slack measures?

% ================================================================
\section{Nine Pre-Declared Empirical Cells}
% ================================================================

Quarterly inflation is considered for three price indexes:
\begin{enumerate}
\item PPI;
\item headline CPI;
\item core CPI.
\end{enumerate}

The activity or marginal-cost block is considered under three measures:
\begin{enumerate}
\item inverse markup;
\item Beveridge--Nelson (BN) output gap;
\item negative unemployment gap.
\end{enumerate}

The resulting design is shown in \cref{tab:nine_cells}.

\begin{table}[htbp]
\centering
\caption{Nine pre-declared empirical cells}
\label{tab:nine_cells}
\begin{tabular}{p{0.36\textwidth}ccc}
\toprule
Inflation / expectation block & Inverse markup & BN output gap & Negative unemployment gap \\
\midrule
PPI / GDP-deflator expectation & Cell 1 & Cell 2 & Cell 3 \\
CPI / CPI expectation & Cell 4 & Cell 5 & Cell 6 \\
Core CPI / CPI expectation & Cell 7 & Cell 8 & Cell 9 \\
\bottomrule
\end{tabular}
\end{table}

The nine cells are frozen before empirical comparison.

\subsection{Tier 1: Quantitative HSA Cell}

\paragraph{Cell 1: PPI \(\times\) inverse markup.}
This is the unique primary quantitative HSA cell.

PPI is the closest observed inflation measure to the producer-pricing equation underlying the HSA restriction, while inverse markup provides the closest direct empirical mapping to real marginal cost.

Only this cell receives the quantitative restriction in \cref{eq:local_hsa_restriction}.

\subsection{Tier 2: Mechanism Robustness}

The following cells move one major dimension away from the theory-near benchmark:
\begin{itemize}
\item Cell 2: PPI \(\times\) BN output gap;
\item Cell 3: PPI \(\times\) negative unemployment gap;
\item Cell 4: CPI \(\times\) inverse markup.
\end{itemize}

These cells ask whether the estimated competition mechanisms survive alternative measures of marginal cost/slack or an alternative inflation index.

They do not receive the quantitative HSA equality.

\subsection{Tier 3: Macro Transportability}

The remaining cells are:
\begin{itemize}
\item Cell 5: CPI \(\times\) BN output gap;
\item Cell 6: CPI \(\times\) negative unemployment gap;
\item Cell 7: core CPI \(\times\) inverse markup;
\item Cell 8: core CPI \(\times\) BN output gap;
\item Cell 9: core CPI \(\times\) negative unemployment gap.
\end{itemize}

These cells study whether the competition mechanisms transport to conventional or underlying consumer-price Phillips curves.

No structural HSA equality is imposed in these cells.

% ================================================================
\section{Inflation and Expectation Measures}
% ================================================================

For price index \(j\in\{PPI,CPI,Core\}\), quarterly annualized inflation is
\begin{equation}
\pi_t^j
=
400
\left(
\log P_t^j-\log P_{t-1}^j
\right).
\label{eq:inflation_general}
\end{equation}

All expectation regressors are one-quarter-ahead forecasts whenever available and are expressed in corresponding quarterly annualized units.

\subsection{PPI Cells}

The preferred conditioning variable is a one-quarter-ahead PPI expectation.

If a sufficiently long PPI expectation series is unavailable, a one-quarter-ahead GDP-deflator expectation is used as an observed conditioning proxy. It is not relabeled as a PPI expectation.

The baseline therefore conditions on
\begin{equation}
g_t^{GDP}
\end{equation}
as a nuisance expectation regressor.

A mandatory robustness exercise introduces an explicit PPI--GDP expectation bridge whenever sufficient information exists to identify it.

\subsection{CPI Cells}

Headline CPI inflation is paired with a one-quarter-ahead CPI expectation.

\subsection{Core CPI Cells}

If a direct core-CPI expectation is unavailable, headline CPI expectation is used as a conditioning proxy. It is explicitly labeled as a headline-CPI expectation proxy rather than an expected core-inflation measure.

A headline-to-core expectation bridge is a robustness exercise when sufficient data are available.

% ================================================================
\section{Marginal-Cost and Slack Measures}
% ================================================================

\subsection{Inverse Markup}

Let
\begin{equation}
\mu_t
\end{equation}
denote the observed markup. Real marginal cost is \(1/\mu_t\).

No sample-mean normalization is used in the structural inverse-markup cells.

Instead, define the marginal-cost deviation relative to the same reference environment used for the HSA elasticity:
\begin{equation}
\boxed{
x_t^{IM}
=
\log(1/\mu_t)
-
\log(1/\mu_{\mathrm{ref}}^f)
=
\log
\left(
\frac{\mu_{\mathrm{ref}}^f}{\mu_t}
\right).
}
\label{eq:inverse_markup_gap}
\end{equation}

Thus \(x_t^{IM}=0\) has an economic interpretation: observed real marginal cost equals its reference flexible-price/steady-state value.

This eliminates dependence on an arbitrary sample-mean centering constant.

\subsection{BN Output Gap}

Let
\begin{equation}
x_t^{BN}
\end{equation}
denote the pre-declared BN output gap.

The zero of the measure already has the interpretation of output being at its estimated permanent component, so no additional centering is applied.

\subsection{Negative Unemployment Gap}

Let
\begin{equation}
u_t^{gap}=u_t-u_t^\ast.
\end{equation}

Define
\begin{equation}
x_t^{U}
=
-u_t^{gap}.
\end{equation}

The sign convention makes larger values correspond to a tighter labor market.

No additional normalization is applied.

The unemployment and BN measures are not interpreted as being in raw log marginal-cost units. Therefore their coefficients are not directly subject to the HSA quantitative equality.

% ================================================================
\section{HSA Theory and Local Restriction}
% ================================================================

Let
\begin{equation}
n=\log N
\end{equation}
denote the theoretical effective number of active firms or varieties.

Under the maintained Rotemberg-HSA producer-pricing environment,
\begin{equation}
\kappa^{th}(n)
=
\frac{\zeta(n)-1}{\chi},
\end{equation}
and
\begin{equation}
\Theta^{th}(n)
=
\frac{1}{\chi}
\frac{1-\rho(n)}{\rho(n)}.
\end{equation}

The HSA identity is
\begin{equation}
\frac{d\kappa^{th}(n)}{dn}
=
\zeta(n)\Theta^{th}(n).
\label{eq:hsa_identity}
\end{equation}

Locally relate the empirical competition coordinate to the theoretical coordinate by
\begin{equation}
q=b_N(n-n_0),
\qquad b_N>0.
\end{equation}

Under a common local rescaling,
\begin{equation}
\frac{d\kappa(q)}{dq}
=
\zeta(q)\theta(q).
\end{equation}

At the reference environment
\begin{equation}
\bar q=0,
\qquad
\hat q=0,
\end{equation}
the primary quantitative HSA restriction is \cref{eq:local_hsa_restriction}.

This is a pointwise local restriction, not a globally exact coefficient function.

It is imposed only in Cell 1.

% ================================================================
\section{Competition Measurement System}
% ================================================================

The latent empirical competition coordinate is
\begin{equation}
q_t=\bar q_t+\hat q_t,
\label{eq:q_decomposition}
\end{equation}
where:
\begin{itemize}
\item \(\bar q_t\) is a slowly moving competitive environment;
\item \(\hat q_t\) is a stationary competition cycle.
\end{itemize}

The state decomposition is a measurement and regularization device rather than an HSA law.

\subsection{Annual Firm Counts}

Let
\begin{equation}
q_y^A
\end{equation}
denote the centered annual firm-count coordinate.

For a point-in-time annual observation at quarter \(t_y\),
\begin{equation}
q_y^A
=
\bar q_{t_y}
+
\hat q_{t_y}
+
\nu_y^N,
\qquad
\nu_y^N
\sim
\Normal(0,\sigma_N^2).
\label{eq:annual_measurement}
\end{equation}

The loading of one defines the scale of the empirical competition coordinate.

The annual firm count is treated as the defining empirical business-demography coordinate, not as an exact observation of theoretical effective varieties.

No additional latent time-varying proxy wedge is included in the baseline unless more than one independently informative annual competition measure is available, because a single annual series cannot separately identify the common competition state and a flexible persistent proxy wedge.

Instead, the design uses:
\begin{itemize}
\item sensitivity to \(\sigma_N\);
\item alternative annual competition proxies;
\item market-universe checks;
\item a multiple-proxy measurement model if multiple credible annual measures become available.
\end{itemize}

\subsection{Quarterly Business-Demography Measurements}

Let
\begin{equation}
e_{jt},
\qquad j=1,\ldots,J_E,
\end{equation}
denote economically comparable quarterly business-demography indicators.

The baseline measurement equation is
\begin{equation}
e_{jt}
=
\bar e_{jt}
+
b_j\hat q_t
+
\xi_{jt}^E,
\qquad
\xi_{jt}^E\sim\Normal(0,\sigma_{E,j}^2),
\label{eq:quarterly_measurement}
\end{equation}
where each series receives its own slow component \(\bar e_{jt}\).

The slow establishment or business-demography components are not tied to \(\bar q_t\).

The loading \(b_j\) is estimated.

Only quarterly indicators that provide genuinely distinct measurement information are included. Additional series are not added merely to increase the number of measurement rows.

If only one credible quarterly establishment series is available, the system reduces to
\begin{equation}
e_t
=
\bar e_t+b_E\hat q_t+\xi_t^E.
\end{equation}

A mandatory conservative sensitivity allows each quarterly indicator to contain an additional idiosyncratic stationary cycle:
\begin{align}
e_{jt}
&=
\bar e_{jt}
+
b_j\hat q_t
+
r_{jt}^E
+
\xi_{jt}^E,\\
r_{jt}^E
&=
\rho_{E,j}r_{j,t-1}^E+\omega_{jt}^E,
\qquad |\rho_{E,j}|<1.
\end{align}

If the common loading or the posterior sharpening of \(\hat q_t\) disappears under this specification, the quarterly data are not considered sufficiently informative about the common competition cycle.

% ================================================================
\section{State Laws}
% ================================================================

The baseline competition-state laws are
\begin{align}
\bar q_t
&=
d_q+\bar q_{t-1}+\eta_t^q,
&
\eta_t^q
&\sim\Normal(0,\sigma_{\bar q}^2),
\\
\hat q_t
&=
\phi_q\hat q_{t-1}+u_t^q,
&u_t^q
&\sim\Normal(0,\sigma_{\hat q}^2),
\qquad |\phi_q|<1.
\end{align}

Each quarterly business-demography slow component obeys a pre-declared smooth law, with a random walk with drift as the baseline:
\begin{equation}
\bar e_{jt}
=
d_{e,j}+\bar e_{j,t-1}+\eta_{jt}^{e}.
\end{equation}

Alternative smooth laws for \(\bar q_t\), \(\hat q_t\), and \(\bar e_{jt}\) are identification sensitivities rather than alternative economic theories.

% ================================================================
\section{Joint Reference-Markup Block}
% ================================================================

The flexible-price/steady-state reference markup is estimated jointly rather than fixed in a preliminary calibration.

Parameterize
\begin{equation}
h_\mu
=
\log(\mu_{\mathrm{ref}}^f-1),
\end{equation}
so that
\begin{equation}
\mu_{\mathrm{ref}}^f
=
1+\exp(h_\mu)>1.
\end{equation}

The associated HSA elasticity is
\begin{equation}
\boxed{
\zeta_0
=
\frac{\mu_{\mathrm{ref}}^f}
{\mu_{\mathrm{ref}}^f-1}
=
1+\exp(-h_\mu).
}
\label{eq:zeta_reference}
\end{equation}

This parameterization guarantees the economically required domain.

A baseline low-frequency markup block is
\begin{equation}
\log\mu_t^{obs}
=
\log\mu_{\mathrm{ref}}^f
+
\lambda_\mu\bar q_t
+
r_t^\mu
+
v_t^\mu,
\qquad
v_t^\mu\sim\Normal(0,\sigma_{\mu,v}^2),
\label{eq:markup_bridge}
\end{equation}
with
\begin{equation}
r_t^\mu
=
\rho_\mu r_{t-1}^\mu+\omega_t^\mu,
\qquad
|\rho_\mu|<1,
\qquad
\omega_t^\mu\sim\Normal(0,\sigma_{\mu,r}^2).
\end{equation}

This equation is an empirical bridge between observed markups and the reference flexible-price markup. It is not treated as an identity equating observed cyclical markups with their flexible-price counterparts.

For every posterior draw,
\begin{equation}
x_t^{IM}
=
\log
\left(
\frac{\mu_{\mathrm{ref}}^f}{\mu_t^{obs}}
\right),
\end{equation}
and, in Cell 1,
\begin{equation}
\kappa_1=\zeta_0\theta_0
\end{equation}
is imposed draw by draw.

Thus reference-markup uncertainty is propagated directly into both the marginal-cost gap and the HSA restriction.

% ================================================================
\section{Inflation Equation}
% ================================================================

For each empirical cell \(c\), the encompassing local equation is
\begin{equation}
\boxed{
\begin{aligned}
\pi_{c,t}
={}&
a_c
+\beta_{b,c}\pi_{c,t-1}
+\beta_{f,c}f_{c,t}^{cond}
+\kappa_{0,c}x_{c,t}
+\kappa_{1,c}\bar q_tx_{c,t}\\
&-\theta_{0,c}\hat q_t
-\gamma_c\bar q_t\hat q_t
+\varepsilon_{c,t}^{\pi},
\end{aligned}
}
\label{eq:inflation_equation}
\end{equation}
with
\begin{equation}
\varepsilon_{c,t}^{\pi}
\sim
\Normal(0,\sigma_{\pi,c}^2)
\end{equation}
in the baseline.

The lagged- and expected-inflation coefficients are estimated separately.

No sum-to-one restriction is imposed.

The coefficient interpretations are
\begin{equation}
\kappa_{1,c}
=
\frac{\partial}{\partial\bar q}
\left(
\frac{\partial\pi_c}{\partial x_c}
\right),
\end{equation}
\begin{equation}
\theta_{0,c}
=
-
\left.
\frac{\partial\pi_c}{\partial\hat q}
\right|_{\bar q=0},
\end{equation}
and
\begin{equation}
\gamma_c
=
-
\frac{\partial^2\pi_c}
{\partial\bar q\,\partial\hat q}.
\end{equation}

% ================================================================
\section{Two-Scale Approximation}
% ================================================================

The baseline is a two-scale varying-coefficient approximation.

The slow state \(\bar q_t\) indexes the competitive environment. Conditional on that environment, \(x_t\) and \(\hat q_t\) are treated as fast cyclical variables.

Thus
\begin{equation}
\kappa(\bar q_t)
\simeq
\kappa_0+\kappa_1\bar q_t,
\end{equation}
and
\begin{equation}
\theta(\bar q_t)
\simeq
\theta_0+\gamma\bar q_t.
\end{equation}

The retained terms are therefore
\begin{equation}
\kappa_0x_t
+\kappa_1\bar q_tx_t
-\theta_0\hat q_t
-\gamma\bar q_t\hat q_t.
\end{equation}

The baseline exclusion of
\begin{equation}
\hat q_tx_t
\end{equation}
is not an HSA zero restriction.

It is the maintained approximation that fast-by-fast interactions are excluded while slow-environment variation in fast-variable coefficients is retained.

A mandatory sensitivity adds
\begin{equation}
\lambda_{qx}\hat q_tx_t
\end{equation}
symmetrically to the relevant restricted and unrestricted models.

A stronger sensitivity imposes
\begin{equation}
\lambda_{qx}=\kappa_1,
\end{equation}
corresponding to a first-order single-coordinate expansion in total
\begin{equation}
q_t=\bar q_t+\hat q_t.
\end{equation}

If the principal result changes materially under either specification, the conclusion is reported as dependent on the two-scale approximation.

% ================================================================
\section{Model Families}
% ================================================================

\subsection{Cell 1: Quantitative HSA Hierarchy}

Cell 1 uses four nested models.

\paragraph{M0: No competition channel.}
\begin{equation}
\kappa_1=\theta_0=\gamma=0.
\end{equation}

\paragraph{M1: Local constant-participation HSA.}
\begin{equation}
\kappa_1=\zeta_0\theta_0,
\qquad
\gamma=0.
\end{equation}

\paragraph{M2: Local HSA.}
\begin{equation}
\kappa_1=\zeta_0\theta_0,
\qquad
\theta_0,\gamma\text{ free}.
\end{equation}

\paragraph{M3: Local unrestricted competition.}
\begin{equation}
\kappa_1,\theta_0,\gamma
\text{ free}.
\end{equation}

The primary HSA residual under M3 is \cref{eq:delta_hsa}.

\subsection{Cells 2--9: Robustness and Transportability Hierarchy}

The HSA equality is not imposed.

Each cell uses:

\paragraph{R0: No competition channel.}
\begin{equation}
\kappa_1=\theta_0=\gamma=0.
\end{equation}

\paragraph{R1: Competition with locally constant participation loading.}
\begin{equation}
\kappa_1,\theta_0
\text{ free},
\qquad
\gamma=0.
\end{equation}

\paragraph{R2: Unrestricted local competition.}
\begin{equation}
\kappa_1,\theta_0,\gamma
\text{ free}.
\end{equation}

Thus the secondary cells answer two clean questions:
\begin{enumerate}
\item do competition-related terms improve the inflation model?;
\item if so, is state dependence of the participation loading empirically required?
\end{enumerate}

No \(\delta_{HSA}\) is interpreted structurally outside Cell 1.

% ================================================================
\section{Full-Joint and Modular-Cut Designs}
% ================================================================

Let \(C\) collect annual and quarterly business-demography measurements.

Let \(S\) collect competition states and their measurement parameters.

For inverse-markup cells, let \(\psi_\mu\) collect reference-markup parameters and latent markup components.

\subsection{Modular Cut}

The primary measurement-supported cut for inverse-markup cells is
\begin{equation}
p_{cut}(S,\psi_\mu,\beta_c)
=
p(S\mid C)
\,p(\psi_\mu\mid \mu^{obs},S)
\,p(\beta_c\mid\pi_c,X_c,S,\psi_\mu).
\label{eq:cut_inverse_markup}
\end{equation}

This ordering has a deliberate interpretation:
\begin{enumerate}
\item business-demography data determine competition states;
\item markup data determine the reference-markup block conditional on those states;
\item inflation determines inflation-equation parameters;
\item inflation does not reweight competition states or the reference-markup calibration.
\end{enumerate}

For BN-output-gap and unemployment-gap cells the middle markup module is absent:
\begin{equation}
p_{cut}(S,\beta_c)
=
p(S\mid C)
p(\beta_c\mid\pi_c,X_c,S).
\label{eq:cut_gap}
\end{equation}

This modular design gives all nine cells the same business-demography-supported competition-state benchmark.

\subsection{Full Joint}

For inverse-markup cells,
\begin{align}
&p_{joint}
(S,\psi_\mu,\beta_c
\mid
C,\mu^{obs},\pi_c,X_c)
\nonumber\\
&\qquad\propto
p(C\mid S)
p(\mu^{obs}\mid S,\psi_\mu)
p(\pi_c\mid S,\psi_\mu,X_c,\beta_c)
p(S)p(\psi_\mu)p(\beta_c).
\label{eq:full_joint_inverse_markup}
\end{align}

For gap cells, the markup block is omitted.

Full-joint inference is a coherent semi-structural pooling estimator.

It is not an endogeneity correction.

Cut--joint differences are treated as information-flow and modular-conflict diagnostics.

% ================================================================
\section{Posterior Simulation}
% ================================================================

The bilinear term
\begin{equation}
-\gamma\bar q_t\hat q_t
\end{equation}
makes the joint state posterior non-Gaussian in the full state vector.

Conditional on either slow or fast state blocks, however, the other block has a linear-Gaussian full conditional.

Posterior inference therefore uses an
\begin{equation}
\boxed{\text{FFBS-within-Gibbs sampler}}
\end{equation}
rather than describing the complete posterior sampler as an exact FFBS algorithm.

Each state-block FFBS step is an exact draw from its corresponding conditional smoothing distribution.

Repeated Gibbs iterations target the joint posterior.

A typical iteration is:
\begin{enumerate}
\item draw \(\hat q_{1:T}\) and the quarterly business-demography slow states conditional on \(\bar q_{1:T}\);
\item draw \(\bar q_{1:T}\) conditional on \(\hat q_{1:T}\);
\item update transition and measurement parameters;
\item update quarterly-measurement loadings and variances;
\item in inverse-markup cells, update the reference-markup block;
\item update inflation-equation parameters;
\item repeat the state blocks before storing a draw if cross-block autocorrelation is high.
\end{enumerate}

Gaussian conditional updates are used where available.

The reference-markup parameter \(h_\mu\), and any other parameter entering the likelihood nonlinearly, is updated with a valid nonconjugate step such as slice sampling or Metropolis--Hastings.

When \(\gamma=0\), a one-block multivariate FFBS implementation is used as a mandatory computational validation benchmark.

% ================================================================
\section{Priors}
% ================================================================

All priors are proper.

Near-improper inverse-gamma variance priors are excluded.

The prior design focuses on:
\begin{itemize}
\item the relative innovation scales of \(\bar q_t\) and \(\hat q_t\);
\item quarterly-measurement trend smoothness;
\item the loading parameters \(b_j\);
\item the annual measurement variance;
\item inflation-equation coefficient scales;
\item the reference-markup parameter \(h_\mu\).
\end{itemize}

The prior on
\begin{equation}
h_\mu=\log(\mu_{\mathrm{ref}}^f-1)
\end{equation}
is based on economically plausible markup magnitudes or external evidence and is not tuned using realized inflation outcomes.

Prior predictive checks examine the implied distributions of
\begin{equation}
\kappa_1\bar q_tx_t,
\qquad
-\theta_0\hat q_t,
\qquad
-\gamma\bar q_t\hat q_t,
\end{equation}
and, in Cell 1,
\begin{equation}
\delta_{HSA}.
\end{equation}

Bayes-factor prior scales are frozen before the model comparisons and are varied over pre-declared economically plausible ranges.

% ================================================================
\section{Identification Diagnostics}
% ================================================================

Identification diagnostics are co-primary with coefficient estimates.

\subsection{Competition-Cycle Measurement}

Compare:

\paragraph{N.}
Annual firms only.

\paragraph{C.}
Annual firms plus quarterly business-demography measurements.

\paragraph{J.}
Full-joint cell-specific model.

\paragraph{J-Q.}
Full-joint model with the quarterly business-demography signal removed.

The central measurement diagnostic is
\begin{equation}
C\text{ versus }N
\end{equation}
for posterior uncertainty about
\begin{equation}
\hat q_t.
\end{equation}

The quarterly data are considered informative only if they materially sharpen the competition cycle and the result survives:
\begin{itemize}
\item alternative quarterly trend laws;
\item alternative competition-cycle laws;
\item idiosyncratic quarterly cycles;
\item alternative quarterly indicators.
\end{itemize}

\subsection{Identification of \(\theta_0\) and \(\gamma\)}

For every main cell report:
\begin{itemize}
\item posterior variation in \(\bar q_t\);
\item posterior variation in \(\hat q_t\);
\item correlation between \(\hat q_t\) and \(\bar q_t\hat q_t\);
\item standardized-design condition number;
\item smallest eigenvalue of the regressor Gram matrix;
\item auxiliary \(R^2\) and VIF for \(\bar q_t\hat q_t\);
\item posterior correlation between \(\theta_0\) and \(\gamma\);
\item prior-to-posterior contraction;
\item cut--joint movement in the coefficients.
\end{itemize}

If \(\gamma\) is not learned from the likelihood, posterior concentration around zero or a Bayes factor favoring \(\gamma=0\) is not interpreted as evidence for local constancy.

\subsection{Persistent Inflation Disturbance}

A mandatory robustness replaces
\begin{equation}
\varepsilon_t^\pi
\sim iid
\end{equation}
with
\begin{equation}
\varepsilon_t^\pi
=
\rho_\pi\varepsilon_{t-1}^\pi+v_t^\pi,
\qquad |\rho_\pi|<1.
\end{equation}

A competition-cycle loading is not considered robust if it collapses once modest residual inflation persistence is admitted.

\subsection{Low-Frequency Inflation Confounding}

A common low-frequency inflation nuisance component is added symmetrically across competing models within each cell.

If \(\kappa_1\) changes materially, the result is described as sensitive to low-frequency inflation specification.

% ================================================================
\section{Measurement-Error and Data-Alignment Checks}
% ================================================================

\subsection{Markup Measurement Error}

Inverse-markup results are conditional on the observed markup unless an identified errors-in-variables structure is available.

A valid measurement-error model requires additional information, such as:
\begin{itemize}
\item externally estimated measurement-error variance;
\item independently constructed markup measures;
\item an external instrument.
\end{itemize}

A measurement equation alone is not treated as identification.

\subsection{Shared Price-Index Construction}

If the markup measure mechanically contains the same price index used to construct inflation, a shared-measurement-error sensitivity is mandatory.

Mechanical covariance is not interpreted as marginal-cost transmission.

\subsection{Market-Universe Alignment}

PPI, CPI/core CPI, markup measures, firm counts, and quarterly business-demography indicators must refer to economically compatible market universes.

Cell 7, core CPI \(\times\) inverse markup, receives particular scrutiny because a broad markup measure may not map cleanly into the core-CPI market universe.

If the alignment is materially poor, the cell is reported as non-interpretable rather than forced into the nine-cell comparison.

% ================================================================
\section{Model Evidence and Predictive Evaluation}
% ================================================================

Model comparison uses two distinct objects.

\subsection{Bayes Factors}

Bayes factors measure integrated model evidence, including fit, parameter uncertainty, prior mass, and model complexity.

They are therefore not described as pure fit measures.

Bayes factors are compared only between models estimated on the same empirical cell.

They are not compared numerically across different dependent variables or different activity measures.

For Cell 1 the main comparisons are:
\begin{equation}
M3\text{ vs }M0,
\end{equation}
\begin{equation}
M2\text{ vs }M3,
\end{equation}
and
\begin{equation}
M1\text{ vs }M2.
\end{equation}

For Cells 2--9 the main comparisons are:
\begin{equation}
R2\text{ vs }R0,
\end{equation}
and
\begin{equation}
R1\text{ vs }R2.
\end{equation}

Formal Bayes factors are based on the full generative model.

No mechanically comparable Bayes factor is assigned to the modular cut distribution.

\subsection{Posterior Predictive Score}

Predictive evaluation uses time-series-respecting posterior predictive densities.

The preferred criterion is recursive one-step-ahead or leave-future-out log predictive score:
\begin{equation}
LPS_M
=
\sum_{t\in\mathcal H}
\log
p(
\pi_t
\mid
\mathcal I_{t-1},M
).
\label{eq:lps}
\end{equation}

Full-sample smoothed latent states are never used to claim forecast performance.

Where contemporaneous activity or markup data are unavailable at the genuine forecast origin, a separate conditional-nowcast exercise is reported.

Predictive-score differences are interpreted within each cell relative to its no-competition benchmark.

\subsection{Posterior Predictive Checks}

Posterior predictive checks are used as absolute-fit diagnostics for:
\begin{itemize}
\item inflation persistence;
\item residual variance;
\item business-demography measurement fit;
\item markup-block fit;
\item the magnitude of competition contributions.
\end{itemize}

Thus the evidence hierarchy is
\begin{equation}
\boxed{
\text{Identification}
\;\rightarrow\;
\text{Bayes factor}
+
\text{predictive score}
\;\rightarrow\;
\text{posterior predictive fit}.
}
\end{equation}

A favorable Bayes factor cannot rescue a model whose key parameter is weakly identified.

% ================================================================
\section{Cross-Cell Interpretation}
% ================================================================

The nine cells are not nine independent attempts to find statistical significance.

They are a fixed robustness architecture.

The interpretation proceeds in the following order.

\paragraph{Gate 1: Competition measurement.}
Does the quarterly business-demography system identify a meaningful \(\hat q_t\) beyond annual firm counts and state-law regularization?

\paragraph{Gate 2: Primary HSA cell.}
In Cell 1:
\begin{itemize}
\item is \(\theta_0\) empirically learned?;
\item is \(\kappa_1\) empirically learned?;
\item is \(\delta_{HSA}\) close to zero?;
\item does Local HSA receive competitive integrated evidence relative to Local U?;
\item do the results survive persistent inflation errors, alternative competition state laws, and the fast-by-fast interaction sensitivity?
\end{itemize}

\paragraph{Gate 3: Marginal-cost/slack robustness.}
Do Cells 2 and 3 produce competition-dependent Phillips-curve slopes and participation loadings with economically compatible signs?

No quantitative HSA equality is required.

\paragraph{Gate 4: Inflation-index transportability.}
Do Cells 4--9 show that the mechanism survives movement from producer inflation to headline and core consumer inflation?

\paragraph{Gate 5: State dependence.}
Only if \(\gamma\) is empirically learned is variation in the participation loading interpreted substantively.

% ================================================================
\section{Required Robustness Exercises}
% ================================================================

The following are mandatory:
\begin{enumerate}
\item annual-firm-only versus quarterly-augmented competition measurement;
\item alternative quarterly business-demography trends;
\item idiosyncratic quarterly business-demography cycles;
\item alternative AR(1)/AR(2) competition-cycle laws;
\item alternative slow competition laws;
\item persistent AR(1) inflation disturbance;
\item low-frequency inflation nuisance component;
\item fast-by-fast interaction \(\hat q_tx_t\);
\item stronger single-coordinate restriction \(\lambda_{qx}=\kappa_1\);
\item markup measurement-error sensitivity where identified;
\item shared-price-index measurement-error check;
\item alternative annual competition proxies;
\item wider annual-measurement variance;
\item expectation-index bridge checks;
\item cut versus full-joint information-flow diagnostics;
\item market-universe alignment;
\item prior sensitivity for all Bayes-factor comparisons;
\item local-approximation range diagnostics.
\end{enumerate}

% ================================================================
\section{Required Simulation Validation}
% ================================================================

Before empirical model ranking, simulation exercises must verify:
\begin{enumerate}
\item agreement between one-block and two-block state samplers when \(\gamma=0\);
\item recovery of \(\theta_0\), \(\gamma\), and \(\kappa_1\) under known parameter values;
\item deterioration of \(\gamma\) identification when variation in \(\bar q_t\) is small;
\item the ability of the C-versus-N diagnostic to reveal weak quarterly information;
\item behavior when quarterly indicators contain idiosyncratic cycles;
\item false loading on \(\hat q_t\) when inflation disturbances are persistent;
\item rejection of the HSA-restricted model when data are generated with economically important \(\delta_{HSA}\neq0\);
\item cut--joint state shifts under a correctly specified joint model;
\item recovery of the reference markup and \(\zeta_0\);
\item RBPF likelihood validation against exact Kalman-filter likelihoods in linear special cases.
\end{enumerate}

% ================================================================
\section{Reporting Plan}
% ================================================================

\subsection{Main Competition-State Figure}

Report
\begin{equation}
\bar q_t,
\qquad
\hat q_t,
\qquad
q_t
\end{equation}
under the business-demography cut and full joint.

Show the C-versus-N uncertainty reduction for \(\hat q_t\).

\subsection{Main Cell-1 Figure}

Report posterior distributions of
\begin{equation}
\kappa_1,
\qquad
\theta_0,
\qquad
\gamma,
\qquad
\delta_{HSA},
\qquad
\mu_{\mathrm{ref}}^f,
\qquad
\zeta_0.
\end{equation}

\subsection{Nine-Cell Coefficient Figure}

Report posterior distributions or credible intervals for
\begin{equation}
\kappa_1,
\qquad
\theta_0,
\qquad
\gamma
\end{equation}
for all nine cells.

The figure visually distinguishes:
\begin{itemize}
\item Tier 1 quantitative HSA;
\item Tier 2 mechanism robustness;
\item Tier 3 transportability.
\end{itemize}

No structural HSA residual is plotted for Cells 2--9.

\subsection{Evidence Table}

For each cell report:
\begin{itemize}
\item Bayes-factor comparisons within that cell;
\item recursive/LFO predictive-score differences relative to the no-competition model;
\item cut and full-joint coefficient summaries;
\item identification diagnostics;
\item persistent-error robustness;
\item convergence and Monte Carlo diagnostics.
\end{itemize}

% ================================================================
\section{Pre-Estimation Freeze}
% ================================================================

Before production estimation, freeze:
\begin{enumerate}
\item all nine empirical cells;
\item Tier 1, Tier 2, and Tier 3 classification;
\item all inflation definitions and forecast horizons;
\item all activity/slack definitions;
\item markup construction;
\item reference-markup model;
\item business-demography sources;
\item annual and quarterly timing conventions;
\item competition-state laws;
\item priors;
\item model families within each cell;
\item modular-cut ordering;
\item full-joint specifications;
\item identification gates;
\item predictive information sets;
\item Bayes-factor prior sensitivity grid;
\item all mandatory robustness exercises.
\end{enumerate}

No cell is added, deleted, or promoted based on the realized coefficient signs or model rankings, except that a cell may be labeled non-interpretable if a pre-declared market-universe or data-quality condition fails.

% ================================================================
\section{Final Interpretation Boundary}
% ================================================================

A favorable result in the primary cell supports the following deliberately conditional statement:

\begin{quote}
Annual firm counts and quarterly business-demography measurements provide an empirically informative latent competition coordinate under the maintained mixed-frequency state-space decomposition. Conditional on the timing, measurement, reference-markup, and shock-orthogonality assumptions, competition changes the marginal-cost slope and is associated with a short-run participation-channel loading. In the PPI/inverse-markup cell, the two channels are locally quantitatively consistent with the HSA restriction after uncertainty in the reference flexible-price markup and the implied demand elasticity is propagated through the joint model.
\end{quote}

The secondary cells establish robustness or transportability only.

The design does not by itself establish:
\begin{itemize}
\item that observed firm counts equal theoretical effective varieties;
\item that the competition state laws are structural;
\item that \(\theta_0\) is a nonparametrically identified causal effect;
\item that \(\gamma\) is a structural pass-through derivative without the additional local proxy-mapping condition;
\item that unemployment or output gaps are in the same units as theoretical marginal cost;
\item that CPI or core CPI obey the producer-pricing HSA equation;
\item that full-joint estimation solves endogenous entry.
\end{itemize}

Those stronger claims require additional measurement structure or independently identified sources of variation.

\end{document}
